\chapter{Phase Histograms}\label{ch:phase-histograms}

This section shows the circular phase histograms per feature for all patients and event types (seizures, \gls{cnn} \glspl{fp}, ensemble \glspl{fp}), based on \cref{sec:cycle-extraction}

The number of events per bin is displayed on the polar y-axis.
The phase angle is displayed around the circle.
The angles indicate the following:
\begin{itemize}
    \item  0°: a phase peak,
    \item  90°: the phase is falling,
    \item 180°: a phase trough,
    \item 270°: the phase is rising.
\end{itemize}

The mean angle denoted after the feature name and also shown as a red arrow.
The relative distance from the plot's center to the arrow's tip represents represents the \gls{plv}.

\foreach \ptnt in \patientList{
    \foreach \eventType in \eventTypes{
        \begin{figure}[!ht]
            \centering
            \includegraphics{phase_histograms/\ptnt/\eventType}
            \captionsetup{list=no}
            \caption{Patient \ptnt\ – \eventType}
            \label{fig:appendix-phase-histogram-\ptnt-\eventType}
        \end{figure}
    }
}

\chapter{Filtered Feature Plots} \label{ch:filtered-feature-plots}
This section contains the filtered features plots per patient and feature.
For each feature, we can see the original feature over time (original signal) and the filtered signal, after using the non-causal filter, as explained in \cref{subsec:cycle-filtering}.
The \glspl{fp} are shown just for the test set, while the seizures are shown for the entire dataset.

\foreach \ptnt in \patientList{
    \section*{Patient \ptnt}
    \foreach \feature in \featureNames{
        \begin{figure}[!ht]
            \centering
            \includegraphics{figures/filtered_feature_plots/\ptnt/\feature}
            \captionsetup{list=no}
            \caption{Patient \ptnt\ -- \feature}
            \label{fig:appendix-filtered-feature-\ptnt-\feature}
        \end{figure}
    }
    \FloatBarrier % This is also necessary so Latex does not register this as a dead loop
}
